% engines: lualatex, pdflatex
% expect-warnings: 0
% expect-errors: 0
\documentclass{article}
\usepackage[margin=1cm]{geometry}
\usepackage{../../latex/pgf-go}
\newcommand\case[1]{\par\medskip\noindent\textbf{#1}\par\nopagebreak}
\newcommand\notes[1][]{\par\noindent Notes: [\pgfgomovenotes[#1]]\par}
\newcommand\ko{\stones[b]{B3,C2,C4}\stones[w]{D2,D4,E3}\moves[w]{C3}\moves{D3}\moves[from = 3, w]{C3}}
\pagestyle{empty}
\begin{document}

\case{Without notes, 3 is drawn on top of 1 / with notes, 3 is not drawn and the notes are "3 at 1"}
\begin{goban}[5]\ko\end{goban}
\begin{goban}[5, notes]\ko\end{goban}
\notes

\case{Several notes, with at = en and separator = \{; \}: "3 en 1; 4 en 2; 5 en 1"}
\begin{goban}[5, notes]\moves{C3, D4}\moves{C3}\moves{D4}\moves{C3}\end{goban}
\notes[at = en, separator = {; }]

\case{Unnumbered moves are not noted: 2 is drawn on 1 (no notes)}
\begin{goban}[5, notes]\moves[numbers = false]{C3}\moves{C3}\end{goban}
\notes

\case{Outside a partial board nothing is noted (no notes); each goban starts without notes}
\begin{goban}[5, notes, partial = A1:B2]\moves{C3, D4}\moves{C3}\end{goban}
\notes

\case{\textbackslash pgfgoset\{notes = \{at = en\}\} and goban = \{notes\} for every board: "2 en 1"}
\begingroup
\pgfgoset{notes = {at = en}, goban = {notes}}
\begin{goban}[5]\moves{C3}\moves[w]{C3}\end{goban}
\notes
\endgroup

\end{document}
